% !TEX root = ../metatime_monograph.tex
\chapter{Octet GMO Coefficients from J-KJ Eigenvalues}
\label{ch:octet_coefficients}

\section{The J-KJ Axial Vector}

The octet mass formula (Eq.~\eqref{eq:gmo_metatime}) has three unknown
coefficients in addition to \(M_0\): \(\alpha\), \(\beta\), and
\(\gamma\).  The coefficient \(\alpha\) is determined directly from the
quark primes; \(\beta\) and \(\gamma\) are derived from \(\alpha\) via
the eigenvalues of the J-KJ (flavour-spin) axial operator.

The J-KJ operator in the Metatime framework is a \(2\times 2\) matrix
in the space spanned by hypercharge \(Y\) and isospin Casimir
\(I(I+1) - Y^2/4\).  Its eigenvalues correspond to the GMO coefficients
\(\beta\) and \(\gamma\).

\section{Coefficient \(\alpha\)}

\begin{equation}
\alpha = E_{\text{proton}}\,\kappa\,
\bigl(q_s^2 - q_u^2 - q_d^2 + L_3\bigr).
\label{eq:alpha}
\end{equation}

Numerically:

\[
\alpha = 938.272 \times 0.00919315 \times (49 - 9 - 25 + 7)
       = 8.62573 \times 22
       = 189.77\ \mev.
\]

\section{Coefficient \(\gamma\)}

The curvature coefficient \(\gamma\) emerges from \(\alpha\) scaled by a
ratio of L-constants:

\begin{equation}
\gamma = \alpha\,\frac{L_4(L_3 - L_4)}{L_3^2}
       = \alpha\,\frac{2(7-2)}{49}
       = \alpha\,\frac{10}{49}.
\label{eq:gamma}
\end{equation}

This form has a clear geometric interpretation:
\begin{itemize}
\item \(L_4(L_3 - L_4) = 2 \times 5 = 10\) is the annihilation--residual
colour mixing.
\item \(L_3^2 = 49\) is the full SU(3) colour-squared phase space.
\item The ratio \(10/49\) is the fraction of \(\alpha\) that contributes
to the SU(3) curvature term.
\end{itemize}

Numerically:

\[
\gamma = 189.77 \times \frac{10}{49}
       = 189.77 \times 0.20408
       = 38.73\ \mev.
\]

\section{Coefficient \(\beta\)}

The linear coefficient \(\beta\) follows from a similar eigenvalue
relation:

\begin{equation}
\beta = -\alpha\,\frac{L_3^2 - 1}{L_3^2}
       = -\alpha\,\frac{48}{49}.
\label{eq:beta}
\end{equation}

The numerator \(L_3^2 - 1 = 49 - 1 = 48\) represents the colour
phase space minus the identity direction (which corresponds to the
singlet \(\Lambda\)).

Numerically:

\[
\beta = -189.77 \times \frac{48}{49}
       = -189.77 \times 0.97959
       = -185.90\ \mev.
\]

\section{Derivation Summary}

\[
\begin{array}{cclc}
\text{Coefficient} & \text{Formula} & \text{Value (MeV)} \\
\hline
\alpha & E_p\,\kappa\,(q_s^2 - q_u^2 - q_d^2 + L_3) & 189.77 \\
\beta   & -\alpha\,(L_3^2 - 1)/L_3^2 & -185.90 \\
\gamma  & \alpha\,L_4(L_3 - L_4)/L_3^2 & 38.73
\end{array}
\]

\section{Comparison with PDG-Fitted Values}

\[
\begin{array}{lccc}
\text{Coefficient} & \text{Metatime} & \text{PDG fit} & \text{Error} \\
\hline
\alpha & 189.77 & 189.73 & +0.02\% \\
\beta   & -185.90 & -183.24 & +1.45\% \\
\gamma  & 38.73 & 38.74 & -0.03\% \\
\end{array}
\]

The \(\gamma\) coefficient is essentially exact.  The \(\beta\) coefficient
has a \(1.45\%\) residual, which propagates into a systematic offset for
the nucleon mass (see Chapter~\ref{ch:octet_predictions}).

\section{Why the Earlier Coefficients Were Incorrect}

The earlier derivation used:

\[
\beta = -\alpha\,\frac{L_3 - 1}{L_3 - 1 + L_4/L_5},
\qquad
\gamma = E_{\text{proton}}\,\kappa\,(L_3 - 1 + L_4/L_5).
\]

These expressions assumed that the J-KJ eigenvalues scale independently
with the L-constants.  The improved derivation shows that \(\gamma\)
scales with \(\alpha\) (not with \(E_{\text{proton}}\kappa\) directly)
and that the correct eigenvalue ratio is \((L_3^2 - 1)/L_3^2\), not
\((L_3-1)/(L_3-1+L_4/L_5)\).
